\section{Related Work}\label{sec:related}

\paragraph{Conformal prediction.}
Split conformal prediction~\cite{vovk2005algorithmic,lei2018distribution} gives marginal coverage under exchangeability, and Mondrian versions calibrate within categories~\cite{vovk2012conditional,ding2023class}. Under covariate shift, weighted conformal prediction uses a known or estimated likelihood ratio~\cite{tibshirani2019conformal,lei2021conformal}, related reweightings handle label shift and general non-exchangeability~\cite{podkopaev2021distribution,barber2023beyond}, and group-weighted conformal prediction (GWCP)~\cite{bhattacharyya2024group} handles shifts in the proportions of known groups; we use GWCP with descriptor bins as a baseline. Localized, robust and multi-source methods reweight by proximity or protect against families of target distributions~\cite{guan2023localized,cauchois2024robust,liu2024multisource}. Conditional calibration gives coverage over a class of covariate shifts by quantile regression with the test point imputed~\cite{gibbs2025conformal}, and related methods give group-conditional coverage~\cite{jung2022batch,bairaktari2025kandinsky}, all for exchangeable points. For grouped data, Dunn et al.~\cite{dunn2023distribution} proposed pooling and subsampling for a new group, Lee et al.~\cite{lee2026hierarchical} proved finite-sample guarantees for pooling (HCP), and later work uses a few observations from the test group~\cite{mallick2026generalized}; Duchi et al.~\cite{duchi2024predictive} calibrate one coverage threshold per environment, and macro-coverage~\cite{bhattacharyya2026macro} averages coverage over classes, in the spirit of our family weighting. In protein science, conformal methods have handled the feedback covariate shift of sequence design~\cite{fannjiang2022conformal} and mining protein databases for function~\cite{boger2025functional}. Table~\ref{tab:guarantees} compares the guarantees: ProtCOVER performs conditional calibration over families with HCP's pooling, and our strongest baseline applies group weighting to the same pooling.

\paragraph{Calibration of variant effect predictors.}
ClinGen calibrated computational scores to evidence strengths for the computational evidence codes PP3 (supporting pathogenicity) and BP4 (supporting benignity)~\cite{pejaver2022calibration}; subsequent work found that genome-wide calibration masks heterogeneous performance across genes~\cite{tejura2024calibration} and proposed gene- and domain-aware calibration~\cite{chen2026gene}, and RaCoon calibrates ESM-1b at the residue level with conditional score distributions~\cite{passi2025racoon}. These methods improve calibration of point probabilities within subgroups seen in training; they do not give distribution-free guarantees, and they condition on gene, domain or residue attributes rather than on homology support relative to the calibration data. Kermut~\cite{groth2024kermut} provides calibrated uncertainty for supervised variant effect regression on deep mutational scans. ProteinGym~\cite{notin2023proteingym} supplies the benchmark and precomputed scores we use, including those of PLM-based and alignment-based predictors~\cite{rives2021biological,frazer2021disease,laine2019gemme,notin2022trancepteve,truong2023poet} and of supervised clinical predictors, which we exclude because they are trained on ClinVar labels~\cite{landrum2018clinvar}. Selective classification~\cite{geifman2017selective} motivates our risk-coverage evaluation of abstention.
